\section*{Chapitre \ref{ch:g2:measure} --- Argent et mesures}

\begin{solution}{exo:g2:measure:1}
Crayon~: cm. Longueur de la classe~: m. Plume~: g. Cartable~: kg.
Bouteille d'eau~: L.
\end{solution}

\begin{solution}{exo:g2:measure:2}
$2$~m $= 200$~cm~; \quad $300$~cm $= 3$~m~; \quad $1$~kg $= 1000$~g.
\end{solution}

\begin{solution}{exo:g2:measure:3}
Par exemple~: $16 = 10 + 5 + 1$~; \ $23 = 20 + 2 + 1$~; \
$35 = 20 + 10 + 5$. (D'autres fa\c{c}ons sont possibles.)
\end{solution}

\begin{solution}{exo:g2:measure:4}
$10 - 6 = 4$~; \quad $20 - 17 = 3$~; \quad $50 - 32 = 18$.
\end{solution}

\begin{solution}{exo:g2:measure:5}
Grande aiguille sur $6$, petite entre $4$ et $5$~: quatre heures et
demie ($4{:}30$). Grande aiguille sur $3$, petite juste apr\`es
$8$~: huit heures et quart ($8{:}15$).
\end{solution}

\begin{solution}{exo:g2:measure:6}
Deux heures et demie~: la grande aiguille pointe en bas vers le $6$,
la petite se situe entre le $2$ et le $3$.
\end{solution}

\begin{solution}{exo:g2:measure:7}
Apr\`es mai~: juin. Avant janvier~: d\'ecembre. (Le mois
d'anniversaire varie.)
\end{solution}

\begin{solution}{exo:g2:measure:8}
D\'epense~: $3 + 4 = 7$. Monnaie~: $10 - 7 = 3$.
\end{solution}

\begin{solution}{exo:g2:measure:9}
$1$~kg $= 1000$~g, ce qui est plus que $800$~g~: le sucre est plus
lourd, de $1000 - 800 = 200$~g.
\end{solution}

\begin{solution}{exo:g2:measure:10}
Par exemple $20 + 10 + 5 + 5 = 40$ et $10 + 10 + 10 + 10 = 40$ (deux
fa\c{c}ons avec exactement quatre pi\`eces ou billets).
\end{solution}
